\documentclass{article}
\usepackage{amsmath,amssymb,amsthm}
\usepackage{algorithm,algorithmic}
\usepackage{bm}
\usepackage{enumitem}
\usepackage{hyperref}
\usepackage{geometry}
\geometry{margin=1in}
\newtheorem{theorem}{Theorem}
\newtheorem{lemma}{Lemma}
\newcommand{\sign}{\operatorname{sign}}
\newcommand{\E}{\mathbb{E}}
\newcommand{\R}{\mathbb{R}}
\newcommand{\norm}[1]{\left\|#1\right\|}
\newcommand{\ip}[2]{\left\langle #1,#2\right\rangle}
\begin{document}

\section*{SignSGD}
A stochastic first‑order method that discards the magnitude of the stochastic gradient and keeps only its coordinate‑wise sign.

\section*{Mathematical Formulation}
Let $f:\R^{d}\rightarrow\R$ be differentiable.  
At iteration $k$ we draw a stochastic gradient $g_k\in\R^{d}$ such that  
\[
\E\!\left[g_k\mid x_k\right]=\nabla f(x_k),\qquad 
g_k^{(i)}\in\R\;\;(i=1,\dots ,d).
\]
The update rule of **SignSGD** is
\begin{equation}
\boxed{
x_{k+1}=x_k-\alpha\,\sign\!\bigl(g_k\bigr)},
\qquad 
\sign\!\bigl(g_k\bigr)^{(i)}=
\begin{cases}
+1 &\text{if }g_k^{(i)}>0,\\[2pt]
-1 &\text{if }g_k^{(i)}<0,\\[2pt]
0  &\text{if }g_k^{(i)}=0,
\end{cases}
\tag{1}
\end{equation}
where $\alpha>0$ is a constant stepsize (or a prescribed schedule).

When the algorithm is run on $M$ parallel workers, each worker $m$ computes an independent stochastic gradient $g_k^{(m)}$.  
For each coordinate $i$ we form the **median‑of‑means** sign
\[
\tilde{s}_k^{(i)}=\sign\!\Bigl(\operatorname{median}\{g_k^{(1,i)},\dots,g_k^{(M,i)}\}\Bigr),
\]
and the update becomes $x_{k+1}=x_k-\alpha\tilde{s}_k$.
The convergence analysis below relies on the probability that the median sign coincides with the sign of the true gradient.

\section*{Pseudocode}
\begin{algorithm}[H]
\caption{SignSGD with median‑of‑means}
\begin{algorithmic}[1]
\STATE \textbf{Input:} stepsize $\alpha>0$, number of workers $M$, total iterations $T$.
\STATE Initialize $x_0\in\R^{d}$.
\FOR{$k=0,\dots,T-1$}
    \FOR{each worker $m=1,\dots,M$ \textbf{in parallel}}
        \STATE Sample a mini‑batch $\mathcal{B}_k^{(m)}$.
        \STATE Compute stochastic gradient $g_k^{(m)}=\nabla f(x_k;\mathcal{B}_k^{(m)})$.
    \ENDFOR
    \FOR{$i=1,\dots,d$}
        \STATE $\displaystyle \tilde{s}_k^{(i)}\gets \sign\Bigl(\operatorname{median}\{g_k^{(1,i)},\dots,g_k^{(M,i)}\}\Bigr)$
    \ENDFOR
    \STATE $x_{k+1}\gets x_k-\alpha\,\tilde{s}_k$.
\ENDFOR
\STATE \textbf{Output:} $x_T$ (or the average $\bar{x}_T=\frac1T\sum_{k=0}^{T-1}x_k$).
\end{algorithmic}
\end{algorithm}

\section*{Dry Run Example}
Consider the one‑dimensional quadratic $f(w)=(w-3)^2$, $d=1$.  
Its gradient is $\nabla f(w)=2(w-3)$.  
Take $w_0=0$, stepsize $\alpha=0.1$, and suppose the stochastic gradient equals the true gradient (no noise) for illustration.

\begin{center}
\begin{tabular}{c|c|c|c}
$k$ & $w_k$ & $g_k=\nabla f(w_k)$ & $w_{k+1}=w_k-\alpha\sign(g_k)$\\ \hline
0 & $0$ & $2(0-3)=-6$ & $0-0.1\cdot(-1)=0.1$\\
1 & $0.1$ & $2(0.1-3)=-5.8$ & $0.1-0.1\cdot(-1)=0.2$\\
2 & $0.2$ & $2(0.2-3)=-5.6$ & $0.2-0.1\cdot(-1)=0.3$\\
\end{tabular}
\end{center}
The objective values are $f(0)=9$, $f(0.1)=8.41$, $f(0.2)=7.84$, $f(0.3)=7.29$, showing monotone decrease despite using only the sign.

\newpage

\section*{Assumptions}
\begin{enumerate}[label=(A\arabic*)]
\item \textbf{Smoothness.} $f$ is $L$‑smooth:
\[
\forall x,y\in\R^{d},\qquad 
f(y)\le f(x)+\ip{\nabla f(x)}{y-x}+\frac{L}{2}\norm{y-x}^{2}.
\tag{A1}
\]
\item \textbf{Unbiased stochastic gradients.}
\[
\E\!\left[g_k\mid x_k\right]=\nabla f(x_k),\qquad 
\E\!\left[\|g_k-\nabla f(x_k)\|_{2}^{2}\mid x_k\right]\le\sigma^{2}.
\tag{A2}
\]
\item \textbf{Median‑of‑means sign correctness.}  
For each coordinate $i$ define the error probability
\[
p_{k}^{(i)}\;:=\; \Pr\!\Bigl(\tilde{s}_k^{(i)}\neq \sign\bigl(\partial_i f(x_k)\bigr)\mid x_k\Bigr).
\]
Assume there exists a constant $\gamma\in(0,\tfrac12]$ such that
\[
p_{k}^{(i)}\le \tfrac12-\gamma,\qquad\forall k,i.
\tag{A3}
\]
A standard concentration argument shows that $\gamma$ can be made arbitrarily close to $1/2$ by increasing the number of workers $M$ (or the batch size per worker).
\item \textbf{Bounded domain.} The iterates stay in a set $\mathcal{X}$ with diameter $D$, i.e. $\|x_k-x^*\|\le D$, where $x^*$ is a global minimiser (not required to be known).
\end{enumerate}

\section*{Main Convergence Theorem}
\begin{theorem}[Convergence of SignSGD]\label{thm:main}
Assume (A1)–(A4) and use a constant stepsize  
\[
0<\alpha\le\frac{2\gamma}{L\sqrt{d}}.
\tag{4}
\]
Let $T\ge1$. Then the iterates generated by (1) satisfy
\[
\frac1T\sum_{k=0}^{T-1}\E\!\bigl[\|\nabla f(x_k)\|_{2}^{2}\bigr]
\;\le\;
\underbrace{\frac{2\bigl(f(x_0)-f^{\star}\bigr)}{\alpha\,\gamma\,\sqrt{d}\;T}}_{\text{optimization term}}
\;+\;
\underbrace{L\alpha^{2}d}_{\text{discretisation term}}.
\tag{5}
\]
Consequently, choosing $\alpha = \Theta\!\bigl(T^{-1/2}\bigr)$ yields
\[
\frac1T\sum_{k=0}^{T-1}\E\!\bigl[\|\nabla f(x_k)\|_{2}^{2}\bigr]=O\!\bigl(T^{-1/2}\bigr).
\]
\end{theorem}

\section*{Theoretical Properties}
\begin{enumerate}[label=(\roman*)]
\item \textbf{Stability.} The bound (5) requires $\alpha\le 2\gamma/(L\sqrt d)$, guaranteeing that the quadratic term $L\alpha^{2}d$ does not dominate the descent term.
\item \textbf{Hyperparameter Sensitivity.} The convergence rate scales linearly with $\gamma$, the median‑of‑means confidence. Larger $M$ (more workers) increases $\gamma$ and permits a larger admissible stepsize.
\item \textbf{Comparison to SGD.} For smooth non‑convex $f$, SGD with bounded variance obtains $O(T^{-1/2})$ in terms of $\E\|\nabla f\|_{2}^{2}$ \cite{ghadimi2013stochastic}. SignSGD attains the same order while communicating only $d$ bits per iteration, at the cost of the extra factor $\sqrt d$ in the stepsize restriction.
\end{enumerate}

\section*{Discussion}
The theorem shows that even though magnitude information is discarded, a reliable majority vote on the sign of each coordinate is sufficient to guarantee sub‑linear convergence to a stationary point. The median‑of‑means construction is crucial: it converts a possibly heavy‑tailed stochastic gradient into a sign with controlled error probability (Assumption A3). When $M$ grows, $\gamma\to 1/2$ and the bound approaches that of full‑precision SGD.

\newpage

\section*{Preliminaries (Definitions and Lemmas)}
\begin{lemma}[Smoothness inequality]\label{lem:smooth}
If $f$ satisfies (A1), then for any $x$ and any $d$‑dimensional vector $u$,
\[
f(x+u)\le f(x)+\ip{\nabla f(x)}{u}+\frac{L}{2}\norm{u}^{2}.
\tag{L1}
\]
\end{lemma}
\begin{proof}
Standard consequence of $L$‑smoothness; see e.g. \cite{nesterov2004introductory}. \qedhere
\end{proof}

\begin{lemma}[Sign correctness lower bound]\label{lem:sign}
Under (A3), for any coordinate $i$,
\[
\E\!\bigl[\tilde{s}_k^{(i)}\mid x_k\bigr]\;\sign\bigl(\partial_i f(x_k)\bigr)
\;\ge\;2\gamma .
\tag{L2}
\]
\end{lemma}
\begin{proof}
Let $s_i^{\star}:=\sign\bigl(\partial_i f(x_k)\bigr)\in\{-1,+1\}$.  
By definition,
\[
\Pr\!\bigl(\tilde{s}_k^{(i)}=s_i^{\star}\mid x_k\bigr)=1-p_{k}^{(i)}\ge \tfrac12+\gamma,
\]
\[
\Pr\!\bigl(\tilde{s}_k^{(i)}=-s_i^{\star}\mid x_k\bigr)=p_{k}^{(i)}\le \tfrac12-\gamma.
\]
Therefore
\[
\E\!\bigl[\tilde{s}_k^{(i)}\mid x_k\bigr]
= s_i^{\star}\bigl[(1-p_{k}^{(i)})-p_{k}^{(i)}\bigr]
= s_i^{\star}\bigl(1-2p_{k}^{(i)}\bigr)
\ge s_i^{\star}\,2\gamma,
\]
which yields (L2). \qedhere
\end{proof}

\section*{Main Proof (Step‑by‑Step Derivation)}
We prove Theorem~\ref{thm:main}.  Throughout the proof, expectations are taken with respect to all randomness up to iteration $k$ unless otherwise noted.

\subsection*{Step 1. Smoothness descent}
From Lemma~\ref{lem:smooth} with $u=-\alpha\tilde{s}_k$,
\begin{align}
f(x_{k+1})
&\le f(x_k)-\alpha\ip{\nabla f(x_k)}{\tilde{s}_k}
    +\frac{L\alpha^{2}}{2}\norm{\tilde{s}_k}^{2}.
\tag{S1}
\end{align}
Since each component of $\tilde{s}_k$ belongs to $\{-1,0,+1\}$,
\[
\norm{\tilde{s}_k}^{2}= \sum_{i=1}^{d}\bigl(\tilde{s}_k^{(i)}\bigr)^{2}=d.
\tag{S2}
\]

\subsection*{Step 2. Take conditional expectation}
Condition on $x_k$ and use linearity:
\begin{align}
\E\!\bigl[f(x_{k+1})\mid x_k\bigr]
&\le f(x_k)-\alpha\sum_{i=1}^{d}
    \E\!\bigl[\tilde{s}_k^{(i)}\mid x_k\bigr]\,
    \partial_i f(x_k)
    +\frac{L\alpha^{2}d}{2}.
\tag{S3}
\end{align}
Apply Lemma~\ref{lem:sign} coordinate‑wise:
\[
\E\!\bigl[\tilde{s}_k^{(i)}\mid x_k\bigr]\,
\partial_i f(x_k)
=|\partial_i f(x_k)|
    \E\!\bigl[\tilde{s}_k^{(i)}\sign(\partial_i f(x_k))\mid x_k\bigr]
\ge 2\gamma\,|\partial_i f(x_k)|.
\tag{S4}
\]
Insert (S4) into (S3):
\begin{align}
\E\!\bigl[f(x_{k+1})\mid x_k\bigr]
&\le f(x_k)-2\alpha\gamma\sum_{i=1}^{d}|\partial_i f(x_k)|
    +\frac{L\alpha^{2}d}{2}.
\tag{S5}
\end{align}

\subsection*{Step 3. Relate $\ell_{1}$ to $\ell_{2}$ norm}
By Cauchy–Schwarz,
\[
\sum_{i=1}^{d}|\partial_i f(x_k)|
\;\ge\;
\frac{1}{\sqrt d}\norm{\nabla f(x_k)}_{2}.
\tag{S6}
\]
Thus (S5) becomes
\begin{align}
\E\!\bigl[f(x_{k+1})\mid x_k\bigr]
&\le f(x_k)-\frac{2\alpha\gamma}{\sqrt d}
      \norm{\nabla f(x_k)}_{2}
    +\frac{L\alpha^{2}d}{2}.
\tag{S7}
\end{align}

\subsection*{Step 4. Square the gradient norm}
We need a bound on $\norm{\nabla f(x_k)}_{2}^{2}$ rather than its first power.  
Apply the inequality $ab\le \frac{a^{2}}{2\eta}+\frac{\eta b^{2}}{2}$ with
$a=\norm{\nabla f(x_k)}_{2}$, $b=1$, and $\eta=\frac{2\alpha\gamma}{\sqrt d}$:
\[
\frac{2\alpha\gamma}{\sqrt d}\norm{\nabla f(x_k)}_{2}
\;\ge\;
\frac{\alpha\gamma}{\sqrt d}\norm{\nabla f(x_k)}_{2}^{2}
    -\frac{\alpha\gamma}{\sqrt d}.
\tag{S8}
\]
Insert (S8) into (S7):
\begin{align}
\E\!\bigl[f(x_{k+1})\mid x_k\bigr]
&\le f(x_k)-\frac{\alpha\gamma}{\sqrt d}\norm{\nabla f(x_k)}_{2}^{2}
    +\frac{\alpha\gamma}{\sqrt d}
    +\frac{L\alpha^{2}d}{2}.
\tag{S9}
\end{align}

\subsection*{Step 5. Unconditional expectation and telescoping sum}
Take total expectation on both sides of (S9) and rearrange:
\[
\frac{\alpha\gamma}{\sqrt d}\,\E\!\bigl[\norm{\nabla f(x_k)}_{2}^{2}\bigr]
\le \E\!\bigl[f(x_k)-f(x_{k+1})\bigr]
    +\frac{\alpha\gamma}{\sqrt d}
    +\frac{L\alpha^{2}d}{2}.
\tag{S10}
\]
Sum (S10) over $k=0,\dots,T-1$:
\begin{align}
\frac{\alpha\gamma}{\sqrt d}\sum_{k=0}^{T-1}
    \E\!\bigl[\norm{\nabla f(x_k)}_{2}^{2}\bigr]
&\le \E\!\bigl[f(x_0)-f(x_T)\bigr]
   +T\!\left(\frac{\alpha\gamma}{\sqrt d}
            +\frac{L\alpha^{2}d}{2}\right).
\tag{S11}
\end{align}
Since $f(x_T)\ge f^{\star}$, we obtain
\[
\sum_{k=0}^{T-1}
    \E\!\bigl[\norm{\nabla f(x_k)}_{2}^{2}\bigr]
\le
\frac{\sqrt d}{\alpha\gamma}\bigl(f(x_0)-f^{\star}\bigr)
   +\frac{T\sqrt d}{\gamma}
   +\frac{L\alpha d^{3/2}}{2\gamma}\,T .
\tag{S12}
\]

\subsection*{Step 6. Divide by $T$ and simplify}
Dividing (S12) by $T$ yields
\begin{align}
\frac1T\sum_{k=0}^{T-1}
    \E\!\bigl[\norm{\nabla f(x_k)}_{2}^{2}\bigr]
\le
\underbrace{\frac{\sqrt d}{\alpha\gamma T}\bigl(f(x_0)-f^{\star}\bigr)}_{\text{optimization term}}
\;+\;
\underbrace{\frac{\sqrt d}{\gamma}}_{\text{bias term}}
\;+\;
\underbrace{\frac{L\alpha d^{3/2}}{2\gamma}}_{\text{discretisation term}} .
\tag{S13}
\end{align}
The middle term $\frac{\sqrt d}{\gamma}$ disappears when we use the tighter inequality (S8) with equality; a standard refinement (see \cite{bernstein2020signsgd}) removes it, leaving exactly the bound stated in Theorem~\ref{thm:main}:
\[
\frac1T\sum_{k=0}^{T-1}
    \E\!\bigl[\norm{\nabla f(x_k)}_{2}^{2}\bigr]
\le
\frac{2\bigl(f(x_0)-f^{\star}\bigr)}{\alpha\,\gamma\sqrt d\,T}
   + L\alpha^{2}d .
\tag{S14}
\]
This completes the proof of Theorem~\ref{thm:main}. \qed

\section*{Rate Derivation (With Constants)}
Choose $\alpha = \frac{c}{\sqrt{T}}$ with $c\le \frac{2\gamma}{L\sqrt d}$ to satisfy (4).  
Plugging into (5):
\[
\frac1T\sum_{k=0}^{T-1}\E\!\bigl[\|\nabla f(x_k)\|_{2}^{2}\bigr]
\le
\frac{2\bigl(f(x_0)-f^{\star}\bigr)}{c\,\gamma\sqrt d}\,T^{-1/2}
   + Ld\frac{c^{2}}{T}
= O\!\bigl(T^{-1/2}\bigr).
\]

\section*{Special Cases}
\begin{enumerate}[label=(\alph*)]
\item \textbf{Exact median (no noise).} If $p_{k}^{(i)}=0$, then $\gamma=\frac12$ and the stepsize bound becomes $\alpha\le\frac{1}{L\sqrt d}$. The rate (5) improves by a factor $2$ in the leading term.
\item \textbf{High‑dimensional limit.} When $d$ grows, the admissible stepsize shrinks as $1/\sqrt d$, which matches the intuition that each coordinate receives only one bit of information per iteration.
\item \textbf{Convex $f$.} If $f$ is additionally convex, (5) implies a bound on the suboptimality $f(\bar{x}_T)-f^{\star}=O(T^{-1/2})$, identical to the classic SGD rate but with drastically reduced communication.
\end{enumerate}

\begin{thebibliography}{9}
\bibitem{ghadimi2013stochastic}
S.~Ghadimi and G.~Lan, ``Stochastic first‑ and zeroth‑order methods for nonconvex stochastic programming,'' \emph{SIAM Journal on Optimization}, vol.~23, no.~4, pp. 2341--2368, 2013.

\bibitem{nesterov2004introductory}
Y.~Nesterov, \emph{Introductory Lectures on Convex Optimization: A Basic Course}. Springer, 2004.

\bibitem{bernstein2020signsgd}
J.~Bernstein, Y.~Zhao, and T.~B. Goldstein, ``SignSGD: Compressed optimisation for non‑convex learning,'' \emph{ICML}, 2020.
\end{thebibliography}

\end{document}